\ProvidesFile{List.tex}[Контейнеры]


\subsection{Контейнеры}

Мы уже выше пользовались контейнером типов и элементов.
Давайте я все же явно проговорю, что это такое.
До появления variadic templates в языке не было способа из коробки хранить списки типов и элементов и их надо было моделировать.
Однако, с появлением variadic templates список типов это просто любой шаблон с произвольным количеством аргументов, которые являются типами.
Вот примеры
\begin{cppcode}
template<class...>
struct List {};

template<class...>
struct MyList {};

using List1 = List<int, double>;
using List2 = List<int, MyList<int, double>, List1>;
\end{cppcode}
Списки элементов разного типа делаются с помощью ключевого слова \verb"auto" так:
\begin{cppcode}
template<auto...>
struct List {};

template<auto...>
struct MyList {};

void func(int) {}

using List1 = List<1, 2, 'c', func>;
using List2 = MyList<1, 2, func>;
\end{cppcode}
Так же можно делать списки с элементами одинакового типа.
Делать специально такой тип имеет смысл только, если вы хотите имплементировать проверку, что элементы действительно являются элементами нужного типа.
В этом случае вам придется запомнить тип:
\begin{cppcode}
template<class T, T...>
struct List {};

using List1 = List<int, 1, 2>;
using List2 = List<char, '1', 'c'>;
\end{cppcode}
Я для различения списков типов и объектов предпочитаю использовать \verb"List" для типов и \verb"Seq" для элементов.

\subsubsection{Список типов и операции}

Мы легко можем создавать разные списки, наполненные элементами.
Так же мы можем сохранять списки в переменных с помощью директивы \verb"using".
\begin{cppcode}
template<class...>
struct List1 {};

using MyList = List1<int, double, char>;
\end{cppcode}
Но теперь хотелось бы иметь в запасе некоторый набор операций со списками, для которого мы хотим иметь удобный синтаксис.
Давайте перечислим эти операции:
\begin{enumerate}
\item Является ли класс списком типов.
К спискам будем относить любой тип вида
\begin{cppcode}
template<class...>
class SomeClass {};

SomeClass<int, double, A, std::vector<B>>; // is a type list
int;                                       // is not a type list
std::vector<int>;                          // is a type list
\end{cppcode}
То есть списком будет являться тип наполненный элементами, а не сам шаблон.
Нам нужна возможность делать такую проверку, чтобы лучше сообщать пользователю об ошибках.
Использование должно выглядеть так:
\begin{cppcode}
isTypeList<List<int, double, A, B>>; // true
isTypeList<int>;                     // false
isTypeList<std::vector<int>>;        // true
\end{cppcode}

\item Проверка типов списков.
А именно, мы можем создать несколько шаблонов для задания списков
\begin{cppcode}
template<class...>
struct List1 {};
template<class...>
struct List2 {};

using MyList1 = List1<int, double>;
using MyList2 = List2<int, double>;
using MyList3 = List2<int, double, char>;
\end{cppcode}
Тут типы списков отличаются, у первого он \verb"List1", а у второго \verb"List2".
Так как мы можем задавать разные виды списков и не фиксируем один специальный, то хотелось бы уметь проверять их на равенство.
\begin{cppcode}
areSame2<MyList1, MyList2>;         // false
areSame<MyList1, MyList2>;          // false
areSame<MyList2, MyList3>;          // true
areSame<MyList1, MyList2, MyList3>; // false
\end{cppcode}

\item Конвертация между разными типами.
Так как мы поддерживаем много разных типов списков, то нужно уметь конвертировать один в другой:
\begin{cppcode}
template<class...>
struct List1 {};

template<class...>
struct List2 {};

using MyList = ConvertTo<List1<int, double>, List2>;
// MyList is List2<int, double>
\end{cppcode}

\item Добавление элементов в конец списка.
\begin{cppcode}
template<class...>
struct List {};

using List1 = List<int, double>;

using List2 = Append<List1, char, double>; // or PushBack
// List2 is List<int, double, char, double>
\end{cppcode}

\item Добавление элементов в начало списка.
\begin{cppcode}
template<class...>
struct List {};

using List1 = List<int, double>;

using List2 = Prepend<char, double, List1>; // or PushFront
// List2 is List<char, double, int, double>
\end{cppcode}

\item Разворачивание списка задом наперед.
\begin{cppcode}
template<class...>
struct List {};

using List1 = List<int, double, char>;

using List2 = Reverse<List1>;
// List2 is List<char, double, int>
\end{cppcode}

\item Доступ к элементу в начале и в конце
\begin{cppcode}
template<class...>
struct List {};

using List1 = List<int, double, char>;

using Type1 = Front<List1>;
// Type1 is int
using Type2 = Back<List1>;
// Type2 is char
\end{cppcode}

\item Удаление элементов в конце и в начале
\begin{cppcode}
template<class...>
struct List {};

using List1 = List<int, double, char>;

using List2 = PopFront<List1>;
// List2 is List<double, char>
using List3 = PopFront<List1, 2>;
// List3 is List<char>
using List4 = PopBack<List1, 2>;
// List4 is List<int>;
using List5 = PopBack<List1, 4>;
// List5 is List<>;
\end{cppcode}

\item Слияние двух и более списков.
Приклеивание списков друг за другом:
\begin{cppcode}
template<class...>
struct List {};

using List1 = List<int, double>;
using List2 = List<char, double>;
using List3 = List<short, double>;
using List4 = Merge<List1, List2, List3>;

// List4 is List<int, double, char, double, short, double>;
\end{cppcode}

\item Запрос размера списка.
\begin{cppcode}
template<class...>
struct List {};

using List1 = List<int, double>;

Size<List1>; // equals 2
\end{cppcode}

\item Проверка на пустоту списка.
\begin{cppcode}
template<class...>
struct List {};

using List1 = List<int, double>;
using List2 = List<>;

isEmpty<List1>; // false
isEmpty<List2>; // true
\end{cppcode}

\item Запросить тип на $i$-ой позиции в списке.
\begin{cppcode}
template<class...>
struct List {};

using List1 = List<int, double, char, short, const int&, const double*>;

using Type = At<List1, 4>;
// Type is const int&
\end{cppcode}

\item Запросить индекс типа в списке (либо индекс первого вхождения, либо список индексов для всех вхождений).
\begin{cppcode}
template<class...>
struct List {};

using List1 = List<int, double, char, double, short, double, const int&, const double*>;

IndexOf1<List1, double>; // equals 1
IndicesOf<List1, double>; // equals Seq<1, 3, 5>
\end{cppcode}
Во втором методе можно дополнительно сообщить тип контейнера, в который надо положить список элементов:
\begin{cppcode}
IndicesOf<List1, double, ListE>; // equals ListE<1, 3, 5>
\end{cppcode}

\item Проверка содержит ли список элемент.
\begin{cppcode}
template<class...>
struct List {};

using List1 = List<int, double, char, double, short, double, const int&, const double*>;

Contains<List1, double>; // true
Contains<List1, const double>; // false
\end{cppcode}

\item Стереть элемент из списка (либо первое вхождение, либо все вхождения).
\begin{cppcode}
template<class...>
struct List {};

using List1 = List<int, double, char, double, short, double, const int&, const double*>;

using List2 = Erase1<List1, double>;
// List2 is List<int, char, double, short, double, const int&, const double*>;
using List3 = EraseAll<List1, double>;
// List3 is List<int, char, short, const int&, const double*>;
\end{cppcode}


\item Избавиться от повторений в списке.
\begin{cppcode}
template<class...>
struct List {};

using List1 = List<int, double, char, double, short, double, const int&, const double*>;

using List2 = NoDuplicates<List1>;
// List2 is List<int, double, char, short, const int&, const double*>;
\end{cppcode}

\item Заменить один элемент на другой (либо первое вхождение, либо все вхождения).
\begin{cppcode}
template<class...>
struct List {};

using List1 = List<int, double, char, double, short, double, const int&, const double*>;

using List2 = Replace1<List1, double, double*>;
// List2 is List<int, double*, char, double, short, double, const int&, const double*>;
using List3 = ReplaceAll<List1, double, double*>;
// List3 is List<int, double*, char, double*, short, double*, const int&, const double*>;
\end{cppcode}

\item Применить процедуру ко всем элементам списка (MAP).
В этом случае, мы считаем, что процедура -- это шаблон вида
\begin{cppcode}
template<class Arg>
using Apply = SomethingDependingOn<Arg>;
\end{cppcode}
Тогда использование MAP выглядит так:
\begin{cppcode}
template<class...>
struct List {};

using List1 = List<int, double, char, double>;

template<class T>
using Vectorize = std::vector<T>;

using List2 = Map<List1, Vectorize>;
//List2 is List<std::vector<int>, std::vector<double>, std::vector<char>, std::vector<double>>
\end{cppcode}

\item Операции \verb"foldl" и \verb"foldl".
Данные операции выполняют бинарную операцию ко всем элементам списка.
Например, объединение большого числа списков, имея операцию только объединения двух, можно сделать так
\begin{cppcode}
template<class...>
struct List {};

using List1 = List<int>;
using List2 = List<double>;
using List3 = List<char>;

using BigList = List<List1, List2, List3>;

using List4 = Foldl<Merge2, List<>, BigList>;
// List4 is List<int, double, char>
\end{cppcode}
Здесь \verb"List<>" выступает в качестве начального элемента.
Математически \verb"foldl" и \verb"foldr" отличаются расстановкой скобок.
На примере объединения множеств $U_1,\ldots, U_k$ с начальным множеством $U_0$, это выглядит так:
\[
\begin{aligned}
\operatorname{foldl}(\cup, U_0, \{U_1,\ldots, U_k\}) &=(\ldots ((U_0 \cup U_1)\cup U_2) \ldots U_{k-1}) \cup U_k)\\
\operatorname{foldr}(\cup, \{U_1,\ldots, U_k\}, U_0) &=
(U_1 \cup (U_2 \ldots (U_{k-1}\cup (U_k \cup U_0))\ldots)
\end{aligned}
\]
Есть разновидность \verb"foldl0" и \verb"foldr0", которым не нужен начальный элемент, но они предполагают, что список аргументов содержит хотя бы один элемент.

\item Операция \verb"zip" для двух списков, возвращающая список из пар.
Тут есть проблема, какой список использовать для пар.
И разрешать ли в качестве аргументов разные типы списков.
Давайте для простоты предположу, что списки съедаются одного типа.
Кроме того, можно эту операцию продолжить на случай более чем двух аргументов.
Тогда \verb"zip" должен возвращать список из кортежей длины равной количеству списков в аргументах.
Как это должно работать:
\begin{cppcode}
template<class...>
struct List {};

using List1 = List<int, double>;
using List2 = List<char, short, int>;
using List3 = List<float, const int&>;

using List4 = Zip<List1, List2>;
// List4 is List<List<int, char>, List<double, short>>;
using List5 = Zip<List1, List2, List3>;
// List5 = List<List<int, char, float>, List<double, short, const int&>>;
\end{cppcode}

\item Отобрать элементы удовлетворяющие условию.
В этом случае предикатом выступает шаблон вида
\begin{cppcode}
template<class Arg>
struct PredicateT {
  static constexpr bool value = /*some logical expression*/;
};
\end{cppcode}
Тогда фильтрация будет выглядеть так
\begin{cppcode}
template<class Arg>
struct isStdVectorT : std::false_type {};
template<class Arg, class Alloc>
struct isStdVectorT<std::vector<Arg, Alloc>> : std::true_type {};

using List1 = List<int, std::vector<double>, char, std::vector<short, MyAllocator<short>>>;

using List2 = Filter<List1, isStdVectorT>;
//List2 is  List<std::vector<double>, std::vector<short, MyAllocator<short>>>
\end{cppcode}

\item Составить список всех подмножеств данного списка.
\begin{cppcode}
template<class...>
struct List {};

using List1 = List<int, double>;

using List2 = PowerSet<List1>;
// List2 is List<List<>, List<double>, List<int>, List<int, double>>;
\end{cppcode}


\item Составить список всех пар элементов данного списка.
\begin{cppcode}
template<class...>
struct List {};

using List1 = List<int, double>;

using List2 = PairSet<List1>;
// List2 is List<List<int, int>, List<int, double>, List<double, int>, List<double, double>>;
\end{cppcode}

\item Составить список всех упорядоченных пар элементов данного списка.
\begin{cppcode}
template<class...>
struct List {};

using List1 = List<int, double>;

using List2 =SymPairSet<List1>;
// List2 is List<List<int, int>, List<int, double>, List<double, double>>;
\end{cppcode}
\end{enumerate}

\subsubsection{Имплементация некоторых из методов}

\paragraph{Проверка на лист типов}

Давайте начнем с простой имплементации, а именно с проверки является ли текущий тип списком типов или нет.
Идея заключается в том, чтобы сделать общий тип для не контейнера и специализацию для контейнера.
В первом случае мы сохраним значение \verb"false", а во втором, значение \verb"true".
Руками это можно сделать так
\begin{cppcode}
template<class T>
struct isTypeListT {
  static constexpr bool value = false;
};

template<template<class...> class List, class... Types>
struct isTypeListT<List<Types...>> {
  static constexpr bool value = true;
};

template<class T>
constexpr bool isTypeList = isTypeListT<T>::value;
\end{cppcode}
Можно немного автоматизировать вычисление \verb"value" следующим образом:
\begin{cppcode}
template<class T>
struct isTypeListT : std::false_type {
};

template<template<class...> class List, class... Types>
struct isTypeListT<List<Types...>> : std::true_type {
};

template<class T>
constexpr bool isTypeList = isTypeListT<T>::value;
\end{cppcode}
Никакой проверки на ошибки тут не надо.

\paragraph{Элемент на $i$-ом месте}

Это один из интересных методов, если мы хотим обрабатывать ошибки.
Простая индуктивная имплементация:
\begin{cppcode}
template<class List, int Index>
struct AtT {
  using type = typename AtT<PopFront<List>, Index - 1>::type;
};

template<class List>
struct AtT<List, 0> {
  using type = Front<List>;
};

template<class List, int Index>
using At = typename AtT<List, Index>::type;
\end{cppcode}
Во-первых мы видим, что очень неудобная рекурсия получается.
Шаблон \verb"AtT" не видит \verb"At", а потому приходится писать длинное выражение с \verb"typename".
Это можно решить так
\begin{cppcode}
template<class List, int Index>
struct AtT;

template<class List, int Index>
using At = typename AtT<List, Index>::type;

template<class List, int Index>
struct AtT {
  using type = At<PopFront<List>, Index - 1>;
};

template<class List>
struct AtT<List, 0> {
  using type = Front<List>;
};
\end{cppcode}
Если у нас реализован \verb"PopFront" для произовольного количества элементов в начале списка, то можно упростить имплементацию еще сильнее:
\begin{cppcode}
template<class List, int Index>
struct AtT;

template<class List, int Index>
using At = typename AtT<List, Index>::type;

template<class List, int Index>
struct AtT {
  using type = Front<PopFront<List, Index>>;
};
\end{cppcode}
Теперь важно добавить обработку ошибок:
\begin{cppcode}
template<class List, int Index>
struct AtT;

template<class List, int Index>
using At = typename AtT<List, Index>::type;

template<class List, int Index>
struct AtT {
  static_assert(isTypeList<List>,
                "The first argument of At MUST be a type List!");
  static_assert(Index >= 0, "The second argument of At MUST be non negative!");
  static_assert(Index < Size<List>,
      "The second argument of At MUST be less the the size of the type List!");
  using type = Front<PopFront<List, Index>>;
};
\end{cppcode}
Тут мы проверяем три условия:
\begin{enumerate}
\item Что класс \verb"List" является действительно списком типов.

\item Что индекс неотрицательный.

\item Что индекс не больше, чем размер списка.
\end{enumerate}
В этом подходе все хорошо, однако у нас в одном классе находятся три \verb"static_assert"-а и еще вызов вспомогательных шаблонов.
В итоге при вызове вроде такого
\begin{cppcode}
using Elem = At<int, -1>;
\end{cppcode}
Вы увидите не только ошибки от двух ассертов из \verb"At", но еще и ошибку из \verb"PopFront".
Это нужно решать специализациями шаблона так, чтобы в проблемных специализациях не было некорректных вызовов других шаблонов.

\paragraph{Избавление от лишнего вывода ошибок}

Для этого можно во вспомогательном шаблоне \verb"AtT" завести флаг на корректность данных.
\begin{cppcode}
template<class List, int Index, bool Level>
struct AtT;

template<class List, int Index>
using At = typename AtT<List, Index, isDataCorrect<List, Index>>::type;

template<class List, int Index>
struct AtT<List, Index, false> {
  static_assert(isTypeList<List>,
                "The first argument of At MUST be a type List!");
  static_assert(Index >= 0, "The second argument of At MUST not be negative!");
  using type = void;
};

template<template<class...> class List, class... Types, int Index>
struct AtT<List<Types...>, Index, false> {
  static_assert(Index >= 0, "The second argument of At MUST not be negative!");
  static_assert(
      Index < Size<List<Types...>>,
      "The second artument of At MUST be less than the size of the type List!");
  using type = void;
};

template<class List, int Index>
struct AtT<List, Index, true> {
  using type = Front<PopFront<List, Index>>;
};
\end{cppcode}
Заметьте, что для обработки ошибок, мы специально делаем две специализации: для списка типов и не для списка типов.
Это связано с тем, что мы не хотим, чтобы функция \verb"Size" бросала ошибки, если \verb"List" не является списком типов.
При этом в обоих случаях надо вставить те ассерты, за счет которых могла быть ошибка.
В первой специализации у нас \verb"List" точно не список, а потому условие на то, что \verb"Index" больше, чем размер списка не имеет смысла и мы его не проверяем.
Во второй специализации у нас наоборот \verb"List<Types...>" всегда список, а потому мы не проверяем условие на список типов.
Зато в этом случае имеет смысл проверить, что \verb"Index" не превосходит размера списка.
Кроме того, чтобы подавить ошибку про отсутствие типа \verb"type" внутри класса, мы вводим тип заглушку вида \verb"void".

\paragraph{Проверка сложного условия}

Теперь разберем как же проверить на корректность данных.
Мы бы могли написать так:
\begin{cppcode}
template<class List, int Index>
constexpr bool isDataCorrect = isTypeList<List> && Index >= 0 && Index < Size<List>;
\end{cppcode}
Однако, при такой имплементации если \verb"List" не является списком типов, то функция \verb"Size" будет кидать ошибки.
А мы этого не хотим.
Потому тут у нас логически ветвление:
\begin{enumerate}
\item Если \verb"List" не список типов, то сразу вернуть \verb"false"

\item Если же \verb"List" список типов, то надо вернуть проверку на индекс.
\end{enumerate}
Делается это все с помощью специализаций.
А для этого надо перейти к классам.
\begin{cppcode}
template<class List, int Index, bool isList>
struct isDataCorrectT : std::false_type {
};

template<class List, int Index>
struct isDataCorrectT<List, Index, true> {
  static constexpr bool value = Index >= 0 && Index < Size<List>;
};

template<class List, int Index>
constexpr bool isDataCorrect = isDataCorrectT<List, Index, isTypeList<List>>::value;
\end{cppcode}
Вот такая имплементация функции \verb"At" и корректности ее данных не будет кидать лишних ошибок и вывод будет намного чище, чем если бы мы пошли прямым путем.

\paragraph{Странная обработка ошибок}

Если вы обратили внимание, то вот тут
\begin{cppcode}
template<class List, int Index>
struct AtT<List, Index, false> {
  static_assert(isTypeList<List>,
                "The first argument of At MUST be a type List!");
  static_assert(Index >= 0, "The second argument of At MUST not be negative!");
  using type = void;
};
\end{cppcode}
Первая проверка выглядит лишней.
Ведь мы зашли в эту специализацию, предполагая, что \verb"List" не является списком типов.
Так почему бы не написать так
\begin{cppcode}
template<class List, int Index>
struct AtT<List, Index, false> {
  static_assert(false,
                "The first argument of At MUST be a type List!");
  static_assert(Index >= 0, "The second argument of At MUST not be negative!");
  using type = void;
};
\end{cppcode}
Однако, мы живем в мире C++, в котором использование \verb"false" напрямую в \verb"static_assert" является UB.
Это супер тупо, но эту проблему приходится решать искусственным образом.
В примере выше, я просто еще раз позвал проверку на список, которая гарантированно даст \verb"false".
Однако, можно сделать вызов шаблона, который всегда даст \verb"false" вне зависимости от аргумента.
Для этого надо ввести вспомогательный условный \verb"false".
Это делается так: мы определяем шаблонную переменную, которая всегда вычисляется в \verb"false".
\begin{cppcode}
template <class...>
struct FalseT : std::bool_constant<false> { };

template <class... Ts>
constexpr bool False = FalseT<Ts...>::value;

static_assert(False<List>, "Error message!");
\end{cppcode}
Есть желание использовать \verb"False<>", но это не сработает, потому что это вычислится явно в \verb"false" из-за отсутствия зависимостей.
И вроде бы это пофиксили в C++23, но я не уверен на 100\%.

\paragraph{Фильтрация элементов}

Сама имплементация этого метода не сложная.
Куда более интересной является имплементация проверки, что аргумент метода является предикатом.
В начале напишем имплементацию метода:
\begin{cppcode}
template<class List, template<class> class Predicate>
struct FilterImpl;

template<class List, template<class> class Predicate>
using Filter = typename FilterImpl<List, Predicate>::type;

template<class List, template<class> class Predicate>
struct FilterImpl {
  static_assert(isTypeList<List>,
                "The first argument of Filter MUST be a type List!");
};

template<template<class...> class List, class... Types, template<class> class Predicate>
struct FilterImpl<List<Types...>, Predicate> {
    static_assert(hasBoolValue<Predicate<Type>>,
                "The second argument of Filter MUST be a predicate!");
  using Tail = Filter<List<Types...>, Predicate>;
  using type =
      IfElse<Predicate<Type>::value, PushFront<Type, Tail>, Tail>;
};

template<template<class...> class List, template<class> class Predicate>
struct FilterImpl<List<>, Predicate> {
  using type = List<>;
};
\end{cppcode}
Теперь давайте попробуем определить проверку, что \verb"Predicate" это предикат.
Мы по-хорошему хотим проверить, что для любого класса \verb"T", класс \verb"Predicate<T>" содержит статическую \verb"constexpr" переменную \verb"value" имеющую тип \verb"bool".
Конечно же для любого типа мы это не проверим.
Потому идея в том, чтобы делать проверку лишь для тех типов, для которых вызывается этот предикат.
В частности мы взываем его для первого типа \verb"Type" в списке.
А потому надо лишь проверить, что \verb"Predicate<Type>" содержит указанную переменную.
Идея проверки держится на простой идее: если в аргументах специализации шаблона происходит ошибка, то такой шаблон просто игнорируется.
Такой принцип называется SFINAE.
То есть мы бы хотели сделать что-то такое:%
\footnote{Имплементация ниже не работает из-за особенностей языка.
Я объясню почему она не работает и приведу работающую.}
\begin{cppcode}
template<class T>
constexpr bool isBool = std::is_same_v<std::remove_cv_t<T>, bool>;

template<class T, bool boolFlag = true>
struct hasBoolValueT : std::false_type {};

template<class T>
struct hasBoolValueT<T, isBool<decltype(T::value)>> : std::true_type {};

template<class T>
constexpr bool hasBoolValue = hasBoolValueT<T>::value;
\end{cppcode}
Такая имплементация не сработает и я объясню почему.
Но она достаточно простая, чтобы понять, что мы хотим и как работает принцип SFINAE.
Заметьте, что у нас есть общий шаблон на строчках~4-5 и специализация на строчках~7-8.
Если набор параметров \verb"T" и \verb"boolValue" подходит под специализацию, то вызывается специализация, если не подходит или ошибка при генерации специализации, то вызывается общий шаблон.
При этом второй аргумент всегда по умолчанию принимает значение \verb"true".

\paragraph{Как это работает}

Какое мы ожидаем поведение от такой имплементации:
\begin{cppcode}
class A {
  static constexpr int value = 0;
};
class B {
  static constexpr bool value = false;
};
class C {
};

std::cout << "has bool value " << hasBoolValue<A> << '\n';
// returns: has bool value 0
std::cout << "has bool value " << hasBoolValue<B> << '\n';
// returns: has bool value 1
std::cout << "has bool value " << hasBoolValue<C> << '\n';
// returns: has bool value 0
\end{cppcode}
Как это работает:
\begin{enumerate}
\item В начале мы пытаемся обратиться к шаблону \verb"hasBoolValueT<T, boolFlag>" с параметрами \verb"A" и \verb"true".

\item Теперь мы должны понять, какой из двух шаблонов лучше подходит для этих параметров.
\begin{enumerate}
\item Первый шаблон \verb"hasBoolValueT<T, boolFlag>".
Под этот шаблон подходит все что угодно, где первый аргумент любой тип \verb"T", а второй -- любое булево значение \verb"boolValue".

\item Второй шаблон \verb"hasBoolValueT<T, isBool<decltype(T::value)>>".
Мы должны подставить вместо \verb"T" класс \verb"A" и посмотреть подходит ли этот шаблон нам или нет.
При вычислении \verb"decltype(A::value)" получаем тип \verb"int".
И потому \verb"isBool" возвращает \verb"false".
Значит получается \verb"hasBoolValueT<A, false>";
Значит вторая специализация для класса \verb"A" подходит только если второй аргумент \verb"false".
\end{enumerate}
Так как мы зовем шаблон \verb"hasBoolValue" с параметрами \verb"A" и \verb"true", то вторая специализация нам не подходит.
Потому будет выбрана первая специализация.
А первая специализация содержит \verb"value" со значением \verb"false".
\end{enumerate}
Теперь посмотрим, что произойдет в случае класса \verb"B".
\begin{enumerate}
\item В начале мы пытаемся обратиться к шаблону \verb"hasBoolValueT<T, boolFlag>" с параметрами \verb"B" и \verb"true".

\item Теперь мы должны понять, какой из двух шаблонов лучше подходит для этих параметров.
\begin{enumerate}
\item Первый шаблон \verb"hasBoolValueT<T, boolFlag>".
Под этот шаблон подходит все что угодно, где первый аргумент любой тип \verb"T", а второй -- любое булево значение \verb"boolValue".

\item Второй шаблон \verb"hasBoolValueT<T,  isBool<decltype(T::value)>>".
Мы должны подставить вместо \verb"T" класс \verb"B" и посмотреть подходит ли этот шаблон нам или нет.
При вычислении \verb"decltype(B::value)" получаем тип \verb"bool".
И потому \verb"isBool" возвращает \verb"true".
Значит получается \verb"hasBoolValueT<B, true>";
Значит вторая специализация для класса \verb"B" тоже подходит для параметров \verb"B" и \verb"true".
\end{enumerate}
Так как вторая специализация более специфичная, то из двух подходящих будет выбрана именно она.
Но вторая специализация содержит \verb"value" со значением \verb"true".
\end{enumerate}
И наконец последний случай -- случай класса \verb"C":
\begin{enumerate}
\item В начале мы пытаемся обратиться к шаблону \verb"hasBoolValueT<T, boolFlag>" с параметрами \verb"C" и \verb"true".

\item Теперь мы должны понять, какой из двух шаблонов лучше подходит для этих параметров.
\begin{enumerate}
\item Первый шаблон \verb"hasBoolValueT<T, boolFlag>".
Под этот шаблон подходит все что угодно, где первый аргумент любой тип \verb"T", а второй -- любое булево значение \verb"boolValue".

\item Второй шаблон \verb"hasBoolValueT<T,  isBool<decltype(T::value)>>".
Мы должны подставить вместо \verb"T" класс \verb"C" и посмотреть подходит ли этот шаблон нам или нет.
При вычислении \verb"C::value" происходит ошибка, так как нет статической переменной \verb"value" внутри класса \verb"C".
Значит этот шаблон игнорируется вовсе в данном случае.
\end{enumerate}
В итоге остается только одна первая специализация, которую компилятор и выбирает.
А первая специализация содержит \verb"value" со значением \verb"false".
\end{enumerate}
В последнем случае важно, что ошибка происходит именно в аргументах шаблона при специализации.

\paragraph{Как заставить это компилироваться}

Как я написал выше, пример который я написал к сожалению не скомпилируется.
Причина в следующем коде:
\begin{cppcode}
template<class T>
struct hasBoolValueT<T, isBool<decltype(T::value)>> : std::true_type {};
\end{cppcode}
Здесь во втором аргументе у нас вычисляется константа, которая зависит от шаблонного параметра \verb"T" из первого аргумента.
В языке запрещено такое вычисление.
Типы могут друг от друга зависеть, а константы нет.
Потому чтобы решить эту проблему, нужно сделать замену \verb"bool" литералов на типы.
Это можно сделать заменой \verb"bool" $\rightarrow$ \verb"std::integral_constant<bool, -->".
Тогда $m\mapsto {}$\verb"std::integral_constant<bool, m>".
И код будет изменен следующим образом:
\begin{cppcode}
template<class T>
constexpr bool isBool = std::is_same_v<std::remove_cv_t<T>, bool>;

template<class T, class Flag = std::true_type>
struct hasBoolValueT : std::false_type {};

template<class T>
struct hasBoolValueT<T, std::integral_constant<bool, isBool<decltype(T::value)>>>
    : std::true_type {};

template<class T>
constexpr bool hasBoolValue = hasBoolValueT<T>::value;
\end{cppcode}
Данное решение делает все, что было описано выше, но теперь это компилируется.
Тут не было содержательных изменений, только особенности языка, чтобы побороться с условиями на аргументы шаблонов и все.
Это, кстати, еще одна из причин, почему люди стараются пользоваться классами, а не значениями.

\paragraph{Еще немного про аргументы}

Есть желание упростить условие внутри второго аргумента и завернуть это во вспомогательный тип.
Это можно сделать разными способами.
Вот несколько примеров, которые могут прийти на ум.
\begin{cppcode}
template<auto V>
constexpr bool hasTypeBool = isBool<decltype(V)>;

template<class T>
struct hasBoolValueT<T, std::integral_constant<bool, hasTypeBool<T::value>>>
    : std::true_type {};
\end{cppcode}
или
\begin{cppcode}
template<class T>
constexpr bool hasValueOfTypeBool = isBool<decltype(T::value)>;

template<class T>
struct hasBoolValueT<T, std::integral_constant<bool, hasValueOfTypeBool<T>>>
    : std::true_type {};
\end{cppcode}
Первый вариант все еще сработает и дает более удобный синтаксис.
Второй же вариант приведет к ошибке компиляции.
Причина тут в том, что во втором случае теперь обращение \verb"T::value" происходит не в аргументах специализации шаблона, а мы просто вызываем шаблон \verb"hasTypeBool" с ошибочным аргументом в строке~2.
То есть при сворачивании выражений в более читаемые, следите за тем, чтобы то, что выдает ошибку, было непосредственно в аргументах специализации.
А все остальное можно свободно прятать во вспомогательную функцию, чтобы сделать аргументы читаемыми.

\paragraph{Имя переменной}

Вот тут у меня только плохие новости.
Механизмов изменять имя переменной в языке нет.
То есть под каждую переменную придется писать очередную проверку снова и снова.
И вот это то самое единственное место, где могут пригодиться макросы.
Надо превратить вот этот код в макрос:
\begin{cppcode}
template<class T, class Flag = std::true_type>
struct hasBoolValueT : std::false_type {};

template<class T>
struct hasBoolValueT<T, std::integral_constant<bool, isBool<decltype(T::value)>>>
    : std::true_type {};

template<class T>
constexpr bool hasBoolValue = hasBoolValueT<T>::value;
\end{cppcode}
Делаем так.
В начале добавим вспомогательный шаблон для проверки типа объекта:%
\footnote{Операция \texttt{\#\#} для препроцессора является конкатенацией строк.}
\begin{cppcode}
template<class T, auto V>
constexpr bool hasTypeAs = std::is_same_v<T, std::remove_cv_t<decltype(V)>>;

#define HAS_MEMBER_OF_TYPE(member, type)                                       \
  template<class T, class Condition = std::true_type>                          \
  struct Has##member##OfType##type##T : std::false_type {};                    \
  template<class T>                                                            \
  struct Has##member##OfType##type##T<                                         \
      T, std::integral_constant<bool, hasTypeAs<type, T::member>>>             \
      : std::true_type {};                                                     \
  template<class T>                                                            \
  constexpr bool has##member##OfType##type =                                   \
      Has##member##OfType##type##T<T>::value;
\end{cppcode}
Теперь надо в начале сгенерировать шаблон для нужного имени и типа:
\begin{cppcode}
HAS_MEMBER_OF_TYPE(x, int)
\end{cppcode}
Этот шаблон развернется в
\begin{cppcode}
  template<class T, class Condition = std::true_type>
  struct HasxOfTypeintT : std::false_type {};
  template<class T>
  struct HasxOfTypeintT<
      T, std::integral_constant<bool, hasTypeAs<int, T::x>>>
      : std::true_type {};
  template<class T>
  constexpr bool hasxOfTypeint = HasxOfTypeintT<T>::value;
\end{cppcode}
И теперь его можно использовать так:
\begin{cppcode}
class A {
public:
  static constexpr int x = 0;
};

class B {
  static constexpr int x = 0;
};

int main() {
  std::cout << "has int x " << hasxOfTypeint<A> << '\n';  // returns true
  std::cout << "has int x " << hasxOfTypeint<B> << '\n';  // returns false
  return 0;
}
\end{cppcode}
Обратите внимание, что тут идет проверка на доступную статическую переменную типа \verb"int" с именем \verb"x" внутри класса.
Правильная проверка на наличие тех или иных полей в классе -- это отдельный вопрос.
% TO DO

\paragraph{Избавление от дубликатов}

При избавлении от дубликатов вопрос в том, кого оставлять.
Стандартная тактика -- оставлять самое первое вхождение или самое последнее.
Давайте имплементируем первый подход.
\begin{cppcode}
template<class List>
struct NoDuplicatesImpl;

template<class List>
using NoDuplicates = typename NoDuplicatesImpl<List>::type;

template<class List>
struct NoDuplicatesImpl {
  static_assert(isTypeContainer<List>,
                "Argument of NoDuplicates MUST be a type List!");
};

template<template<class...> class List, class Type1, class... Types>
struct NoDuplicatesImpl<List<Type1, Types...>> {
private:
  using GoodTail = NoDuplicates<List<Types...>>;
  using TailWithoutHead = Erase1<GoodTail, Type1>;

public:
  using type = Prepend1<Type1, TailWithoutHead>;
};
\end{cppcode}
Если же надо оставить самые последние вхождения, то можно имплементировать так
\begin{cppcode}
template<template<class...> class List, class Type1, class... Types>
struct NoDuplicatesImpl<List<Type1, Types...>> {
private:
  using GoodTail = NoDuplicates<List<Types...>>;

public:
  using type = IfElse<Contains<Type1, GoodTail>, GoodTail, Prepend<Type1, GoodTail>>;
};
\end{cppcode}

\paragraph{Множество подмножеств}

Это одна из интересных операций, которая легко делается с помощью \verb"MAP" и рекурсии.
\begin{cppcode}
template<class List>
struct PowerSetImpl;

template<class List>
using PowerSet = typename PowerSetImpl<List>::type;

template<class List>
struct PowerSetImpl {
  static_assert(isTypeList<List>,
                "Argument of PowerSet MUST be a type List!");
};

template<template<class...> class List, class Type, class... Types>
struct PowerSetImpl<List<Type, Types...>> {
  using TailPowerSet = typename PowerSet<List<Types...>>;

  template<class T>
  using AddHead = PushFront<Type, T>;

  using HeadWithOthers = Map<TailPowerSet, AddHead>;

  using type = Merge2<HeadWithOthers, TailPowerSet>;
};

template<template<class...> class List>
struct PowerSetImpl<List<>> {
  using type = List<List<>>;
};
\end{cppcode}
Логика в этой имплементации следующая.
Если есть список из типов $L = \{t_1,\ldots, t_n\}$.
То в начале составим все подмножеств для $S = 2^{\{t_2, \ldots, t_n\}}$.
После чего надо к каждому множеству из $S$ добавить пропущенный $t_1$ и получим $S'$.
Тогда $S\cup S'$ и будет множеством всех подмножеств в $L$.

\paragraph{Множество пар}

Еще одна интересная операция -- список всех пар.
\begin{cppcode}
template<class List>
struct PairSetImpl;

template<class List>
using PairSet = typename PairSetImpl<List>::type;

template<class List>
struct PairSetImpl {
  static_assert(isTypeList<List>,
                "Argument of PairSet MUST be a type List!");
};

template<template<class...> class List, class Type, class... Types>
struct PairSetImpl<List<Type, Types...>> {
  using TailPairs = PairSet<List<Types...>>;

  template<class T>
  using AddHeadFront = List<Type, T>;

  using HeadListPairs = Map<List<Type, Types...>, AddHeadFront>;

  template<class T>
  using AddHeadBack = List<T, Type>;

  using TailHeadPairs = Map<List<Types...>, AddHeadBack>;

  using type = Merge<HeadListPairs, TailHeadPairs, TailPairs>;
};

template<template<class...> class List>
struct PairSetImpl<List<>> {
  using type = List<>;
};
\end{cppcode}
Визуально этот подход можно представлять себе так.
Если у нас есть список типов $L = \{t_1,\ldots, t_n\}$, то все пары можно изобразить в следующем квадрате
\[
\xymatrix@R=10pt@C=10pt{
  {}&{t_1}&{t_2}&{\ldots}&{t_n}\\
  {t_1}&{\phantom{A}2}{\save
   [].[rrr]*+[F-]\frm{}
  \restore}&{}&{}&{}\\
  {t_2}&{\phantom{A}3}{\save
   [].[dd]*+[F-]\frm{}
  \restore}&{\phantom{A}1}{\save
   [].[rrdd]*+[F-]\frm{}
  \restore}&{}&{}\\
  {\vdots}&{}&{}&{}&{}\\
  {t_n}&{}&{}&{}&{}\\
}
\]
Здесь соответствие следующее:
\begin{enumerate}
\item \verb"TailPairs".
Это просто все пары в списке без первого элемента.
В этом месте идет рекурсия.

\item \verb"HeadListPairs".
Это мы строим пары $(t_1, t)$, где $t$ пробегает весь список $L$ (включая $t_1$).

\item \verb"TailHeadPairs".
Это мы строим все пары вида $(t, t_1)$, где $t$ пробегает хвост, то есть $L \setminus\{t_1\}$.
\end{enumerate}
И в конце все три списка сливаются в один в порядке $2$, $3$, $1$.

\subsection{Список элементов}

Список элементов разного типа -- это любой шаблон вида
\begin{cppcode}
template<auto...>
struct List {};
\end{cppcode}
Тогда в нем можно хранить элементы следующим образом:
\begin{cppcode}
struct A {
  void f() const {
  }
  static void g() {
  }
};

void f() {
}

template<auto...>
struct List {};

using MyList = List<1, 'c', f, &A::f, A::g>;
\end{cppcode}
Обратите внимание, что обычные функции и статические функции можно складывать без взятия указателя с помощью \verb"&".

\subsubsection{Список операций}

Список операций тут такой же как и у списка типов, но с минимальными изменениями.
Из интересного:
\begin{enumerate}
\item Можно добавить операцию \verb"TypeOf", которая возвращает тип $i$-ого элемента.

\item Так как операцию над элементами можно делать двумя способами: с помощью шаблонов и с помощью \verb"constexpr" функций, то методы принимающие функции (например \verb"Map" или \verb"Foldl") могут принимать метод двумя разными способами.
К сожалению или к счастью нельзя иметь шаблон с разными аргументами, потому приходится делать два шаблона.
\end{enumerate}

\subsubsection{Имплементация некоторых}

Давайте начнем с \verb"Map".
Начнем с того, как можно использовать этот метод.
Предположим, что у нас есть список целых чисел и мы хотим каждый элемент возвести в квадрат.
\begin{cppcode}
template<auto...>
struct List {};

using List1 = List<1, 2, 3, 4>;
\end{cppcode}
Метод возводящий элементы в квадрат можно имплементировать двумя способами.
\begin{enumerate}
\item С помощью шаблона
\begin{cppcode}
template<int arg>
struct SquareT {
  static constexpr int value = arg * arg;
};

using List2 = MapT<List1, SquareT>;
// List 2 is List<1, 4, 9, 16>;
\end{cppcode}

\item С помощью \verb"constexpr" метода
\begin{cppcode}
constexpr int square(int arg) {
  return arg * arg;
}

using List2 = Map<List1, square>;
// List2 is List<1, 2, 9, 16>
\end{cppcode}
\end{enumerate}
Давайте покажем как может выглядеть имплементация подобных методов.
В начале шаблонная имплементация.
\begin{cppcode}
template<class List, template<auto> class F>
struct MapTImpl {
  static_assert(isElemList<List>,
                "The first argument of Map MUST be an element list!");
};

template<template<auto> class F, template<auto...> class List,
         auto... Elements>
struct MapTImpl<List<Elements...>, F> {
  // static_assert(F is a function, "The second argument of Map MUST be a function!");
  using type = List<F<Elements>::value...>;
};

template<class List, template<auto> class F>
using MapT = typename MapTImpl<List, F>::type;
\end{cppcode}
Теперь имплементация для \verb"constexpr" функций.
\begin{cppcode}

template<class List, auto F>
struct MapImpl {
  static_assert(isElemList<List>,
                "The first argument of Map MUST be an element list!");
};

template<template<auto...> class List, auto... Elements, auto F>
struct MapImpl<List<Elements...>, F> {
  // static_assert(F is a function, "The second argument of Map MUST be a function!");
  using type = List<F(Elements)...>;
};

template<class List, auto F>
using Map = typename MapImpl<List, F>::type;
\end{cppcode}
Проверку, что \verb"F" является функцией, я убрал.
Ее можно сделать по аналогии с тем, как это было разобрано для случая \verb"Filter" для списка типов.

\paragraph{Что еще}

Еще полезно иметь возможность конвертации списка элементов в типы.
Это всегда делается с помощью шаблонов, так как \verb"constexpr" функции не могут вернуть тип, только значение.
И в целом можно наворотить хитрую проверку, если шаблон метода содержит \verb"value", то значит делаем список данных, а если шаблон метода содержит тип \verb"type", то конвертируем к списку типов.

\subsection{Список однотипных элементов}

Если мы хотим работать со списком однотипных элементов, то у нас появляется возможность проверять тип элементов, прежде чем мы их положим в список.
Так же можно сравнивать типы при слиянии списков и прочих операциях.
Это дает дополнительную защиту от ошибок программиста.
А потому обработка ошибок -- это главное, чем этот случай отличается от общего.
Вот как может выглядеть список элементов одного типа:
\begin{cppcode}
template<class T, T... Elements>
struct List {};

using List1 = List<int, 1, 2, 3>;
using List2 = List<char, 'a', 'b', 'c'>;
\end{cppcode}

\subsubsection{Имплементация некоторых методов}

Вот как может выглядеть добавление одного элемента в начало списка с проверкой типов.
\begin{cppcode}
template<auto Element, class List>
constexpr bool hasTypeAsInList = std::is_same_v<decltype(Element), TypeOf<List>;

template<auto Element, class List>
struct PushFront1Impl {
  static_assert(isElemList<List>,
                "The second argument of PushFront1 MUST be an element List!");
  static_assert(hasTypeAsInList<Element, List>,
                "The type of the first argument of PushFront1 MUST be the same as "
                "the type of the elements in the List!");
};

template<class Type, template<class, Type...> class List, Type Element,
         Type... Elements>
struct PushFront1Impl<Element, List<Type, Elements...>> {
  using type = List<Type, Element, Elements...>;
};

template<auto Element, class List>
using PushFront1 = typename PushFront1Impl<Element, List>::type;
\end{cppcode}
Слияние двух списков можно имплементировать так:
\begin{cppcode}
template<class T1, class T2>
struct Merge2Impl {
  static_assert(isElemList<T1>,
                "First argument of Merge MUST be an element container!");
  static_assert(isElemList<T2>,
                "Second argument of Merge MUST be an element container!");
  static_assert(areSame2<T1, T2>, "First and Second arguments of Merge MUST be "
                                  "the same element containers!");
  static_assert(std::is_same_v<TypeOf<T1>, TypeOf<T2>>,
                "Types of elements in containers in Merge MUST be the same!");
};

template<class Type, template<class, Type...> class List,
         Type... Elements1, Type... Elements2>
struct Merge2Impl<List<Type, Elements1...>,
                  List<Type, Elements2...>> {
  using type = List<Type, Elements1..., Elements2...>;
};

template<class T1, class T2>
using Merge2 = typename Merge2Impl<T1, T2>::type;
\end{cppcode}
Мы обычно не думаем про проверку ошибок, когда пишем код.
Но это в целом плохая привычка.
И кроме того, для данного списка проверка ошибок -- это одна из основных фич, которые нужно имплементировать.
